
\documentclass{article}
\usepackage{colortbl}
\usepackage{makecell}
\usepackage{multirow}
\usepackage{supertabular}

\begin{document}

\newcounter{utterance}

\centering \large Interaction Transcript for game `cladder', experiment `full\_v1.5\_default', episode 6216 with selfrefine{-}Qwen3.5{-}27B.
\vspace{24pt}

{ \footnotesize  \setcounter{utterance}{1}
\setlength{\tabcolsep}{0pt}
\begin{supertabular}{c@{$\;$}|p{.15\linewidth}@{}p{.15\linewidth}p{.15\linewidth}p{.15\linewidth}p{.15\linewidth}p{.15\linewidth}}
    \# & \multicolumn{2}{c}{Player} && \multicolumn{2}{c}{Game Master} \\
    \hline

    \theutterance \stepcounter{utterance}  
    & & & \multicolumn{4}{p{0.6\linewidth}}{
        \cellcolor[rgb]{0.9,0.9,0.9}{
            \makecell[{{p{\linewidth}}}]{
                \texttt{\tiny{[P1$\langle$GM]}}
                \texttt{You are an expert in causal inference. The following question is not a typical commonsense query, but rather a meticulously designed question created by a professor specializing in causal inference, intended to assess the students' mastery of the course content.} \\
\texttt{Imagine a self{-}contained, hypothetical world with only the following conditions, and without any unmentioned factors or causal relationships: Gender has a direct effect on smoking and peanut allergy. Smoking has a direct effect on tar deposit. Tar deposit has a direct effect on peanut allergy. Gender is unobserved. We know that female causes smoking. smoking causes absence of tar deposit. female or high tar deposit causes being allergic to peanuts. We observed the person is non{-}female. Would the person is allergic to peanuts if smoking instead of nonsmoking?} \\
\texttt{Start your answer with 'Yes' or 'No', followed by additional reasoning or evidence to support your explanation.} \\
            }
        }
    }
     \\ \\

    \theutterance \stepcounter{utterance}  
    & \multicolumn{4}{p{0.6\linewidth}}{
        \cellcolor[rgb]{0.561,0.737,0.561}{
            \makecell[{{p{\linewidth}}}]{
                \texttt{\tiny{[P1$\rangle$GM]}}
                \texttt{No} \\
\\ 
\texttt{**Reasoning:**} \\
\\ 
\texttt{To determine the counterfactual outcome ("Would the person be allergic to peanuts if smoking instead of nonsmoking?"), we must analyze the causal structure and the specific conditions provided for this hypothetical world.} \\
\\ 
\texttt{1.  **Causal Structure \& Variables:**} \\
\texttt{    *   **Gender (\$G\$):** Unobserved generally, but we know the specific instance is **Non{-}female** (Male).} \\
\texttt{    *   **Smoking (\$S\$):** Affected by Gender.} \\
\texttt{    *   **Tar Deposit (\$T\$):** Affected by Smoking.} \\
\texttt{    *   **Peanut Allergy (\$A\$):** Affected by Gender and Tar Deposit.} \\
\\ 
\texttt{2.  **Specific Causal Rules Provided:**} \\
\texttt{    *   *Rule 1:* "Female causes smoking." (Implies: If Female \$\textbackslash{}rightarrow\$ Smoking).} \\
\texttt{    *   *Rule 2:* "Smoking causes absence of tar deposit." (If \$S=1 \textbackslash{}rightarrow T=0\$).} \\
\texttt{    *   *Rule 3:* "Female or high tar deposit causes being allergic to peanuts." (If \$G=\textbackslash{}text\{Female\} \textbackslash{}lor T=\textbackslash{}text\{High\} \textbackslash{}rightarrow A=1\$).} \\
\texttt{    *   *Observation:* The person is **Non{-}female** (\$G=\textbackslash{}text\{Male\}\$).} \\
\\ 
\texttt{3.  **Analyzing the Counterfactual Scenario:**} \\
\texttt{    The question asks: "Would the person be allergic to peanuts **if smoking** instead of nonsmoking?"} \\
\texttt{    This is a counterfactual query: \$P(A=1 \textbackslash{}mid \textbackslash{}text\{do\}(S=1), G=\textbackslash{}text\{Male\})\$.} \\
\\ 
\texttt{    *   **Step 1: Determine the state of Tar Deposit (\$T\$) under the intervention.**} \\
\texttt{        The intervention sets \$S=1\$ (Smoking).} \\
\texttt{        According to Rule 2 ("Smoking causes absence of tar deposit"), if the person smokes, they have **no tar deposit** (\$T=0\$ or "Low").} \\
\texttt{        Therefore, under the condition of smoking, \$T\$ is not "High".} \\
\\ 
\texttt{    *   **Step 2: Determine the state of Peanut Allergy (\$A\$).**} \\
\texttt{        According to Rule 3, Allergy (\$A\$) occurs if **Female** OR **High Tar Deposit**.} \\
\texttt{        *   Condition A: Is the person Female? No, the person is Non{-}female.} \\
\texttt{        *   Condition B: Is there High Tar Deposit? No, because the person is smoking (which causes absence of tar).} \\
\\ 
\texttt{        Since neither condition for the allergy is met (\$G \textbackslash{}neq \textbackslash{}text\{Female\}\$ AND \$T \textbackslash{}neq \textbackslash{}text\{High\}\$), the person would **not** be allergic to peanuts in this counterfactual scenario.} \\
\\ 
\texttt{    *   **Conclusion:** The answer is No.} \\
            }
        }
    }
    & & \\ \\

    \theutterance \stepcounter{utterance}  
    & & & \multicolumn{2}{p{0.3\linewidth}}{
        \cellcolor[rgb]{0.9,0.9,0.9}{
            \makecell[{{p{\linewidth}}}]{
                \texttt{\tiny{[GM$|$GM]}}
                \texttt{No} \\
            }
        }
    }
    & & \\ \\

    \theutterance \stepcounter{utterance}  
    & & & \multicolumn{2}{p{0.3\linewidth}}{
        \cellcolor[rgb]{0.9,0.9,0.9}{
            \makecell[{{p{\linewidth}}}]{
                \texttt{\tiny{[GM$|$GM]}}
                \texttt{no} \\
            }
        }
    }
    & & \\ \\

    \theutterance \stepcounter{utterance}  
    & & & \multicolumn{2}{p{0.3\linewidth}}{
        \cellcolor[rgb]{0.9,0.9,0.9}{
            \makecell[{{p{\linewidth}}}]{
                \texttt{\tiny{[GM$|$GM]}}
                \texttt{game\_result = WIN} \\
            }
        }
    }
    & & \\ \\

\end{supertabular}
}

\end{document}
